% !TEX root = ../notes_dgxspark.tex
% Chapter 1: Quickstart -- from your laptop to a shell on the Spark
%%%%%%%%%%%%%%%%%%%%%%%%%%%%%%%%%%%%%%%%%%%%%%%%%%%%%%%%%%%%%%%%%%%%%%
%
% Access path: WireGuard tunnel into the lab network, then ssh. Accounts
% and tunnel configs are issued by Mark by mail. mDNS does not cross the
% tunnel, so the chapter uses the IP address from that mail, never the
% spark-xxxx.local name.

\index{quickstart}\index{DGX Spark}
\chapter{Quickstart.}
\printchapterversion{01_quickstart}

\section{Ignition}

\index{account}\newthought{The Spark has one user account per person.}%
\marginnote{\newthought{This setup is far from perfect.} If you want to
build a job scheduler, per-user VPN configurations, or sandboxes that keep
users apart, reach out.}
Request yours by mail. The reply carries your username and
password, and the setup folder as an attached zip with the WireGuard
configuration and the Spark's IP address.
\marginnote{\newthought{Scan the code} and send a mail stating your name, why you need access, which
course, and what you want to run.\par\medskip
\null\hfill\qrcode[height=1.0cm]{mailto:schutera@dhbw-ravensburg.de?subject=DGX\%20Spark\%20Account\&body=Name\%3A\%20\%0ACourse\%3A\%20\%0AProject\%3A\%20}}

\index{sudo}\index{shared infrastructure}\newthought{Your account is an
admin account} on a shared machine. With \texttt{sudo} you can install what
you need, and also break the machine for everyone. With great power comes great
responsibility: Store nothing sensitive on it, leave other users' files and
processes alone, sandbox your stuff and never touch the system config itself. 
Protect yourself from others, and protect others from yourself. 
While you have your own user, two users being logged in at the same time
will produce lag and slowdowns. If you experience this, please log out and try again later.
\marginnote{\newthought{Remember?} Chapter 1 of
Full Stack Handwerkszeug, Development and Production, covers this split,
along with SSH and Linux servers.
\url{https://material.schutera.com/lecturenotes/notes_missingsemester/notes_missingsemester.pdf}}%
That also means the Spark is used most efficiently when you spend as little time as possible to deploy your jobs and leave, 
not for interactive development. If you need to develop interactively, please do so on your own machine.
\marginnote{\newthought{You learn useful stuff here. qed.}}  


\marginnote{\newthought{Download.} \url{https://www.wireguard.com/install/}}
\index{WireGuard}\newthought{WireGuard opens a tunnel into the lab
network.} Only traffic to the lab goes through it; the rest of your
connection stays as it is. Install the WireGuard client, choose
\texttt{Import tunnel(s) from file}, pick \texttt{dgx-spark.conf} from the
setup folder, and press \texttt{Activate}.
The status should now read
\texttt{Active} and \texttt{Latest handshake} is seconds old. No handshake
means your network blocks the tunnel. 

\index{SSH}\newthought{Then log in} with the IP address from the mail:
\begin{verbatim}
ssh <username>@<spark-ip>
\end{verbatim}

\index{nvidia-smi}\newthought{Before you start a job,} run
\texttt{nvidia-smi} to see who is on the GPU.

